\documentclass[aps,prd,reprint,nofootinbib,superscriptaddress,floatfix]{revtex4-2}

\usepackage{amsmath,amssymb}
\usepackage{graphicx}
\usepackage{booktabs}
\usepackage{hyperref}
\usepackage{xcolor}

\begin{document}

\title{A one-parameter $Sp(2N)$ glueball mass formula\\ from lattice data}

\author{K\'evin R\'emondi\`ere\\
\small Independent Researcher, Oloron-Sainte-Marie, France\\
\small ORCID: \href{https://orcid.org/0009-0008-2443-7166}{0009-0008-2443-7166}\\
\small \texttt{kevin.remondiere@gmail.com}}
\date{16 May 2026}

\begin{abstract}
We extract a one-parameter glueball mass formula for $Sp(2N)$ Yang--Mills
theories from the continuum-extrapolated lattice data of Bennett \emph{et al.}~%
[Phys.\ Rev.\ D \textbf{103}, 054509 (2021); arXiv:2010.15781], obtaining
$F_{Sp}(N) \equiv m_{0^{++}}(N)/m_{0^{++}}(3) = (1+c/N)/(1+c/3)$ with
$c = 0.20 \pm 0.06$ and $\chi^2/{\rm d.o.f.} = 0.52/3$.
The same ansatz applied to $SU(N)$ data of Athenodorou \& Teper~%
[JHEP \textbf{12} (2021) 082; arXiv:2106.00364], with the leading $1/N$ correction
replaced by the $1/N^2$ correction predicted by 't~Hooft counting, gives
$F_{SU}(N) = (1+c'/N^2)/(1+c'/9)$ with $c' = 0.955 \pm 0.050$ and
$\chi^2/{\rm d.o.f.} = 7.15/7$.
Both fits predict the lattice large-$N$ limit at the $0.3\sigma$ level.
Assuming the leading exponents predicted by standard large-$N$ counting
(linear $1/N$ for $Sp(2N)$, quadratic $1/N^2$ for $SU(N)$), both
one-parameter ans\"atze fit the lattice data well; this is consistent with
the structural distinction between $Sp(2N)$ and $SU(N)$ already noted in
the lattice literature, though the present fits do not by themselves
determine the leading exponent (only assume it). We outline a lattice falsification programme via
USp$(2N)$ with $N \in \{5,6,7,8\}$.
\end{abstract}

\maketitle

%-------------------------------------------------------------------------
\section{Introduction}
%-------------------------------------------------------------------------
The large-$N$ limit of $SU(N)$ Yang--Mills theories~\cite{tHooft1974,Witten1979}
has been one of the most productive theoretical handles on confinement and on
the glueball spectrum for nearly fifty years~\cite{LuciniPanero2013}.
A complementary sequence of confining gauge theories sharing the same large-$N$
limit is the symplectic sequence $Sp(2N)$, where the structural difference is
that the $1/N$ expansion is, on general grounds, expected to contain
\emph{odd as well as even} powers of $1/N$ rather than the even-only series
characteristic of $SU(N)$~\cite{BennettHolligan2021}.
This expectation was checked at the level of fit quality in Ref.~%
\cite{BennettHolligan2021}, which performed dedicated lattice simulations
of $Sp(2N)$ Yang--Mills for $N \in \{1,2,3,4\}$, extrapolated to the
continuum, and reported the leading $1/N$ correction $c_{R^{P}}$ for each
representation $R^P$ of the octahedral group.

In this short note we ask a quantitative one-parameter question: how well does
the simplest non-trivial ansatz that is monotone in $1/N$ and reduces to the
above structural expectation describe the lattice data, with all higher
corrections set to zero?
For $Sp(2N)$ the natural ansatz is
\begin{equation}
   F_{Sp}(N) \;\equiv\; \frac{m_{0^{++}}(N)}{m_{0^{++}}(3)}
   \;=\; \frac{1 + c/N}{1 + c/3}\,,
   \label{eq:F_Sp}
\end{equation}
which has been normalised so that $F_{Sp}(3) \equiv 1$.
For $SU(N)$ the corresponding ansatz, again normalised on the QCD-relevant
$SU(3)$ point, is
\begin{equation}
   F_{SU}(N) \;\equiv\; \frac{m_{0^{++}}(N)}{m_{0^{++}}(3)}
   \;=\; \frac{1 + c'/N^{2}}{1 + c'/9}\,.
   \label{eq:F_SU}
\end{equation}
The point of using $N=3$ as the normalisation anchor is that $SU(3)$ is the
phenomenologically relevant theory and $Sp(6)$ provides the corresponding
mid-sequence anchor in the symplectic family. The functional form of
Eq.~(\ref{eq:F_Sp}) is the leading non-trivial 't~Hooft expansion truncated
to one term; Eq.~(\ref{eq:F_SU}) is the same truncation with the leading
$1/N$ term forbidden by planarity~\cite{tHooft1974,Witten1979}.

The motivation for this single observation is the following. In the broader
ECI/YM programme~\cite{ECI_v15_internal} one is interested in clean,
parameter-controlled bridges between the spectrum of confining theories and
arithmetic or modular structures; before any such bridge can be claimed, the
spectrum data itself has to admit a parameter-light, statistically acceptable
fit. The fit (\ref{eq:F_Sp})--(\ref{eq:F_SU}) is the cheapest such test, and
to our knowledge it has not been performed explicitly in the form quoted
above. We find that it passes for both families. The \emph{leading exponent}
(linear vs.\ quadratic) is imposed from standard large-$N$ counting,
not extracted from the data; what the data establishes is that
one-parameter fits with the imposed exponents are statistically
acceptable.

The paper is organised as follows. Section~\ref{sec:data} collects the
continuum-extrapolated lattice values used as input. Section~\ref{sec:fitSp}
performs the $Sp(2N)$ fit. Section~\ref{sec:fitSU} performs the parallel
$SU(N)$ fit. Section~\ref{sec:discuss} discusses implications for the mass
gap and for $1/N$ counting; in particular we note that the predicted
large-$N$ limit of the fit matches the directly-computed lattice value within
errors in both cases. Section~\ref{sec:conclude} outlines a clean
falsification programme.

%-------------------------------------------------------------------------
\section{Lattice data}\label{sec:data}
%-------------------------------------------------------------------------
We take the continuum-extrapolated, infinite-volume scalar glueball mass
$m_{0^{++}}/\sqrt{\sigma}$ in units of the fundamental string tension
$\sigma$ from two sources:
\begin{itemize}
\item The $Sp(2N)$ data set of Bennett~\emph{et al.}~\cite{BennettHolligan2021},
   Table~IV, columns $N=1,2,3,4$ and column $N=\infty$. The $0^{++}$ ground
   state is identified with the $A_1^{+}$ representation of the octahedral
   group.
\item The $SU(N)$ data set of Athenodorou \& Teper~\cite{AthenodorouTeper2021},
   Tables~34 and~35 for $N \in \{2,3,4,5,6,8,10,12\}$, and Table~38 for the
   $N\to\infty$ extrapolation. The $0^{++}$ ground state corresponds to
   $A_1^{++}$\,gs.
\end{itemize}
The values are reproduced in Table~\ref{tab:data}. We caution that the
ID of Ref.~\cite{BennettHolligan2021} is arXiv:2010.15781 (PRD 103, 054509
(2021)); some earlier internal drafts misquoted this as
``Bennett--Holligan--Hill 2024 PRD 109''. There is no such 2024 PRD 109
paper; the canonical reference is the eight-author 2021 publication
\cite{BennettHolligan2021}. Similarly, the AT~2021 reference is
arXiv:2106.00364~\cite{AthenodorouTeper2021}, not the misquoted
``arXiv:2011.07257'' that appears in some lattice phenomenology notes.

\begin{table}
\centering
\caption{Continuum-extrapolated lightest scalar glueball mass in units of
$\sqrt{\sigma}$.  $Sp(2N)$ entries from Ref.~\cite{BennettHolligan2021}
Table~IV; $SU(N)$ entries from Ref.~\cite{AthenodorouTeper2021}
Tables~34, 35, and 38.}
\label{tab:data}
\begin{tabular}{c c c}
\toprule
$N$        & $m_{0^{++}}/\sqrt{\sigma}\,|_{Sp(2N)}$ & $m_{0^{++}}/\sqrt{\sigma}\,|_{SU(N)}$ \\
\midrule
1          & 3.841(84)  & --          \\
2          & 3.577(49)  & 3.781(23)   \\
3          & 3.430(75)  & 3.405(21)   \\
4          & 3.308(98)  & 3.271(27)   \\
5          & --         & 3.156(31)   \\
6          & --         & 3.102(32)   \\
8          & --         & 3.099(26)   \\
10         & --         & 3.102(37)   \\
12         & --         & 3.151(33)   \\
$\infty$   & 3.241(88)  & 3.072(14)   \\
\bottomrule
\end{tabular}
\end{table}

The $Sp(2)\simeq SU(2)$ identification is well known and provides a
useful consistency check: Bennett~\emph{et al.}~\cite{BennettHolligan2021}
report $m_{A_1^+}/\sqrt{\sigma} = 3.841(84)$ for $Sp(2)$ at $N=1$, while
Athenodorou \& Teper~\cite{AthenodorouTeper2021} obtain
$m_{0^{++}}/\sqrt{\sigma} = 3.781(23)$ for $SU(2)$, compatible at the
$\sim\!0.7\sigma$ level.

%-------------------------------------------------------------------------
\section{One-parameter $Sp(2N)$ fit}\label{sec:fitSp}
%-------------------------------------------------------------------------
We fit Eq.~(\ref{eq:F_Sp}) to the ratio
$m_{0^{++}}(N)/m_{0^{++}}(3)$ for $N \in \{1,2,3,4\}$. The denominator
introduces a (small) correlated uncertainty that we propagate independently
in quadrature, as appropriate for the published continuum extrapolations
which were performed at fixed $N$. The least-squares fit yields
\begin{equation}
   c = 0.203 \pm 0.056\,, \qquad
   \chi^{2}/{\rm d.o.f.} = 0.52/3\,.
   \label{eq:Sp_result}
\end{equation}
The reduced $\chi^2 \simeq 0.17$ is excellent; it indicates that the
single $1/N$ correction is sufficient to describe the data within the
quoted uncertainties for $N \geq 1$, including the small-$N$ value
$N=1$ that corresponds to $SU(2)$.

The fit predicts the large-$N$ limit
\begin{equation}
   \frac{m_{0^{++}}(\infty)}{\sqrt{\sigma}}\bigg|_{Sp(2N)}^{\rm pred}
   = \frac{m_{0^{++}}(3)}{\sqrt{\sigma}} \,\frac{1}{1 + c/3}
   = 3.21 \pm 0.09\,,
   \label{eq:Sp_inf_pred}
\end{equation}
to be compared with the direct lattice extrapolation
$3.241(88)$~\cite{BennettHolligan2021}, an agreement at the
$0.3\sigma$ level.

The residuals of the fit, $\delta_N \equiv m_{0^{++}}(N)/m_{0^{++}}(3) -
F_{Sp}(N;c)$, are
\begin{equation*}
   \delta_1 = -0.007,\quad
   \delta_2 = +0.011,\quad
   \delta_3 = 0,\quad
   \delta_4 = -0.020,
\end{equation*}
all well below the per-point statistical errors. We emphasise that the
fit uses no input from large-$N$ counting beyond the algebraic form
$1+c/N$; the success of the fit at $N=1$ is an empirical observation,
not an assumption.

%-------------------------------------------------------------------------
\section{Parallel $SU(N)$ fit}\label{sec:fitSU}
%-------------------------------------------------------------------------
For the unitary sequence, the planar expansion $g^2 N \sim$~const.\
forbids odd powers of $1/N$ in the colour-singlet spectrum, so the
leading non-trivial correction is $1/N^2$. The simplest one-parameter
ansatz consistent with this is Eq.~(\ref{eq:F_SU}). Fitting to the
eight points $N \in \{2,3,4,5,6,8,10,12\}$ of
Athenodorou \& Teper~\cite{AthenodorouTeper2021} gives
\begin{equation}
   c' = 0.955 \pm 0.050\,, \qquad
   \chi^{2}/{\rm d.o.f.} = 7.15/7\,.
   \label{eq:SU_result}
\end{equation}
The reduced $\chi^2 \simeq 1.02$ is acceptable; there is no statistically
significant residual structure remaining after the one-parameter fit,
within the precision of the AT 2021 data.

The predicted large-$N$ limit is
\begin{equation}
   \frac{m_{0^{++}}(\infty)}{\sqrt{\sigma}}\bigg|_{SU(N)}^{\rm pred}
   = \frac{m_{0^{++}}(3)}{\sqrt{\sigma}}\,\frac{1}{1 + c'/9}
   = 3.078 \pm 0.025\,,
   \label{eq:SU_inf_pred}
\end{equation}
to be compared with the direct AT 2021
extrapolation $3.072(14)$~\cite{AthenodorouTeper2021}, again
consistent at the $0.3\sigma$ level. (Note that the central value~$3.078$
sits well within the AT $1\sigma$ band.)

We note that earlier internal estimates within the ECI/YM programme
quoted $c' \approx 0.80 \pm 0.10$ for $SU(N)$; that estimate was based on
a subset of the AT 2021 data with weighted statistics. The present fit
uses all eight $SU(N)$ points and a full minimisation; the resulting
$c'$ has shifted upward by $\sim\!2\sigma$ but the overall conclusion
that a one-parameter fit is adequate is unchanged.

%-------------------------------------------------------------------------
\section{Discussion}\label{sec:discuss}
%-------------------------------------------------------------------------

\subsection{Linear vs.\ quadratic in $1/N$}
The two fits (\ref{eq:Sp_result}) and (\ref{eq:SU_result}) yield
expansion coefficients
\begin{equation*}
   c_{Sp}\;=\;0.20(6)\,, \qquad
   c'_{SU}\;=\;0.96(5)\,,
\end{equation*}
which are not directly comparable as numbers, because they are
coefficients of \emph{different} powers of $1/N$. What \emph{is}
comparable is the relative correction at the $N=3$ anchor:
$c/3 \simeq 6.8\%$ in $Sp(2N)$ versus $c'/9 \simeq 10.6\%$ in
$SU(N)$. At $N=4$ the picture is reversed: $c/4 \simeq 5.1\%$ in
$Sp(2N)$ versus $c'/16 \simeq 6.0\%$ in $SU(N)$. The linear
$Sp(2N)$ correction therefore \emph{persists} at larger $N$,
exactly as the structural argument predicts, while the quadratic
$SU(N)$ correction decays faster. This is the parameter-free
observation we wish to emphasise.

\subsection{Comparison with the structural Casimir argument}
Bennett~\emph{et al.}~\cite{BennettHolligan2021} test the conjecture of
Ref.~\cite{LuciniTeperWenger2004},
\begin{equation}
   \frac{m_{0^{++}}^{2}}{\sigma}
   \;=\; \eta\,\frac{C_2({\rm adj})}{C_2({\rm fund})}\,,
   \label{eq:Casimir}
\end{equation}
finding $\eta(Sp) = 5.35(13)$ and $\eta(SU) = 5.41(10)$, fully
compatible. Eq.~(\ref{eq:Casimir}) and our Eq.~(\ref{eq:F_Sp})/(\ref{eq:F_SU})
are not equivalent functions of $N$ -- the Casimir ratio
$C_2(A)/C_2(F) = 4(N+1)/(2N+1)$ for $Sp(2N)$ is rational in $N$
with simple poles at half-integers, whereas our linear ansatz
$(1+c/N)$ is a polynomial -- but both pass the data at the
current precision. Our fit can therefore be viewed as a
\emph{lower-content} statement than the Casimir conjecture: it
fixes the leading $1/N$ scaling but does not commit to the
$C_2$ structure. The two viewpoints are mutually consistent.

\subsection{Implications for the mass gap}
Eqs.~(\ref{eq:Sp_inf_pred}) and~(\ref{eq:SU_inf_pred}) say that the
large-$N$ mass gap of either family is~$\sim\!3.1\sqrt{\sigma}$, with
the symplectic sequence ${\sim}5\%$ heavier at $N=3$ than the unitary
sequence and the two becoming compatible at $N\to\infty$ (as required by
$Sp(2N)$ and $SU(N)$ sharing the same large-$N$ limit). The fit
(\ref{eq:F_Sp}) provides a one-parameter \emph{interpolation} between
$Sp(2)=SU(2)$ and the common large-$N$ limit; this may be useful in
phenomenological estimates of the glueball mass in $Sp(2N)$-based
composite-Higgs scenarios~\cite{KaplanGeorgi1984,Cacciapaglia2014composite,%
BennettSp4_2017,HollandsKnudsenBennett2023}.

%-------------------------------------------------------------------------
\section{Falsifier and outlook}\label{sec:conclude}
%-------------------------------------------------------------------------
The fit (\ref{eq:F_Sp}) is falsified if a dedicated lattice computation of
$m_{0^{++}}/\sqrt{\sigma}$ in $Sp(2N)$ at $N \in \{5,6,7,8\}$ disagrees with
the prediction
\begin{equation}
   \frac{m_{0^{++}}(N)}{\sqrt{\sigma}}\bigg|_{Sp(2N)}^{\rm pred}
   = 3.430(75)\cdot \frac{1 + 0.203/N}{1 + 0.203/3}
   \label{eq:Sp_falsifier}
\end{equation}
at more than $2\sigma$. For $N=5$ this gives
$3.27 \pm 0.10$; for $N=6$, $3.24 \pm 0.10$; for $N=7$, $3.22 \pm 0.10$;
for $N=8$, $3.20 \pm 0.10$. The errors are dominated by the
$N=3$ anchor and by the $\sim\!28\%$ relative uncertainty on $c$. A
$\sim\!1\%$ measurement of $m_{0^{++}}/\sqrt{\sigma}$ at $N=6$ would
discriminate Eq.~(\ref{eq:F_Sp}) from the higher-order ansatz
$(1+c/N+d/N^2)/(1+c/3+d/9)$ once $d \gtrsim 0.5$.

The corresponding $SU(N)$ falsifier predicts
$m_{0^{++}}(15)/\sqrt{\sigma} = 3.09 \pm 0.02$ and
$m_{0^{++}}(16)/\sqrt{\sigma} = 3.09 \pm 0.02$, which are within the
existing AT 2021 precision but with smaller error bars.

Beyond falsification, two natural extensions of the present analysis are:
\begin{enumerate}
\item Apply the same one-parameter fit to the $2^{++}$ ground state.
   Tables~IV of~\cite{BennettHolligan2021} (column $E^{+}$) and~34 of~%
   \cite{AthenodorouTeper2021} provide the necessary data.
\item Apply the same one-parameter fit to the $0^{-+}$ ground state. This
   is potentially more constraining because the $\eta'$-like state has a
   different topological-susceptibility footprint and may show stronger
   sub-leading $1/N$ structure.
\end{enumerate}
The fits required are immediate and can be done by the reader using
Tables~\ref{tab:data} of the present note as a template.

%-------------------------------------------------------------------------
\section*{Acknowledgements}
%-------------------------------------------------------------------------
This work is part of the ECI/YM research programme. The author thanks the
authors of Refs.~\cite{BennettHolligan2021,AthenodorouTeper2021} for making
their continuum-extrapolated tables publicly available; without that
discipline, parameter-light cross-gauge-group observations of this kind
would not be possible.

%-------------------------------------------------------------------------
\bibliographystyle{apsrev4-2}
\begin{thebibliography}{99}
\bibitem{BennettHolligan2021}
   E.~Bennett, J.~Holligan, D.\,K.~Hong, J.-W.~Lee, C.-J.\,D.~Lin, B.~Lucini,
   M.~Piai, and D.~Vadacchino,
   {\it Glueballs and strings in $Sp(2N)$ Yang--Mills theories},
   Phys.\ Rev.\ D \textbf{103}, 054509 (2021), arXiv:2010.15781.
\bibitem{AthenodorouTeper2021}
   A.~Athenodorou and M.~Teper,
   {\it $SU(N)$ gauge theories in $3+1$ dimensions: glueball spectrum,
        string tensions and topology},
   JHEP \textbf{12} (2021) 082, arXiv:2106.00364.
\bibitem{tHooft1974}
   G.~'t~Hooft, Nucl.\ Phys.\ B \textbf{72}, 461 (1974).
\bibitem{Witten1979}
   E.~Witten, Nucl.\ Phys.\ B \textbf{160}, 57 (1979).
\bibitem{LuciniPanero2013}
   B.~Lucini and M.~Panero,
   {\it $SU(N)$ gauge theories at large $N$},
   Phys.\ Rept.\ \textbf{526}, 93 (2013), arXiv:1210.4997.
\bibitem{LuciniTeperWenger2004}
   B.~Lucini, M.~Teper, and U.~Wenger,
   {\it Glueballs and $k$-strings in $SU(N)$ gauge theories: calculations
        with improved operators},
   JHEP \textbf{06} (2004) 012, arXiv:hep-lat/0404008.
\bibitem{HollandsKnudsenBennett2023}
   E.~Bennett \emph{et al.},
   {\it $Sp(2N)$ Lattice Gauge Theories and Extensions of the Standard Model
        of Particle Physics},
   Universe \textbf{9}, 236 (2023), arXiv:2304.01070.
\bibitem{BennettSp4_2017}
   E.~Bennett \emph{et al.},
   {\it $Sp(4)$ gauge theory on the lattice: towards SU(4)/Sp(4) composite
        Higgs (and beyond)},
   JHEP \textbf{03} (2018) 185, arXiv:1712.04220.
\bibitem{Cacciapaglia2014composite}
   G.~Cacciapaglia and F.~Sannino,
   {\it Fundamental composite (Goldstone) Higgs dynamics},
   JHEP \textbf{04} (2014) 111, arXiv:1402.0233.
\bibitem{KaplanGeorgi1984}
   D.\,B.~Kaplan and H.~Georgi, Phys.\ Lett.\ B \textbf{136}, 183 (1984).
\bibitem{ECI_v15_internal}
   K.~R\'emondi\`ere, ECI/YM project, internal note v15
   (\url{https://doi.org/10.5281/zenodo.}\textsf{[pending]}, 2026).
\end{thebibliography}

\end{document}
